\section{Correction Events}

The \texttt{canvass-correction} fixture models a small but important case. The
unofficial snapshot reports one fewer vote for Candidate B in precinct P-001
than the canvassed snapshot. The package includes both sets of summaries and a
public correction event from unofficial to canvassed status.

\begin{center}
\small
\begin{tabularx}{\textwidth}{lrrr}
\toprule
Snapshot & Candidate A & Candidate B & Counted ballots \\
\midrule
P-001 unofficial & 40 & 34 & 79 \\
P-001 canvassed & 40 & 35 & 80 \\
Jurisdiction unofficial & 65 & 64 & 139 \\
Jurisdiction canvassed & 65 & 65 & 140 \\
\bottomrule
\end{tabularx}
\end{center}

The verifier checks three things. First, each summary satisfies local contest
arithmetic: candidate votes plus residual counts equal counted ballots. Second,
each status-specific jurisdiction total equals the sum of its reporting units.
Third, a correction event exists with authority, source references, explanation,
and a transition from unofficial to canvassed.

This does not prove that the correction was lawful. It proves that the package
does not hide the delta. A public reader can see the before snapshot, the after
snapshot, the event metadata, and the source hashes that bind the package to its
declared evidence.

